% 57e Olympiade Internationale de Mathématiques — Hong Kong, Chine.
% Énoncés officiels en français (imo-official.org). Chaque problème vaut 7 points.

\begin{center}
  \emph{Premier jour --- lundi 11 juillet 2016 --- Durée : 4 heures 30}
\end{center}

\subsection*{Problème 1}
Le triangle $BCF$ est rectangle en $B$. Soit $A$ le point de la droite $(CF)$ tel que $FA = FB$ et que $F$ se
trouve entre $A$ et $C$. Le point $D$ est choisi tel que $DA = DC$ et que la droite $(AC)$ soit la bissectrice
de l'angle $\widehat{DAB}$. Le point $E$ est choisi tel que $EA = ED$ et que $(AD)$ soit la bissectrice de
l'angle $\widehat{EAC}$. Soit $M$ le milieu du segment $[CF]$, et $X$ le point tel que $AMXE$ soit un
parallélogramme (c'est-à-dire $(AM) \parallel (EX)$ et $(MX) \parallel (EA)$).

Montrer que les droites $(BD)$, $(FX)$ et $(ME)$ sont concourantes.

\subsection*{Problème 2}
Trouver tous les entiers strictement positifs $n$ pour lesquels les cases d'un tableau $n \times n$ puissent
chacune être remplies avec l'une des lettres $I$, $M$ et $O$ de sorte que les deux conditions suivantes soient
satisfaites :
\begin{itemize}
  \item dans chaque ligne et chaque colonne, exactement un tiers des cases contient la lettre $I$, un tiers
  contient la lettre $M$ et un tiers contient la lettre $O$ ;
  \item dans chaque diagonale composée d'un nombre de cases divisible par trois, exactement un tiers des cases
  contient la lettre $I$, un tiers contient la lettre $M$ et un tiers contient la lettre $O$.
\end{itemize}

\textbf{Note} : les lignes et les colonnes d'un tableau $n \times n$ sont numérotées de $1$ à $n$. Ainsi
chaque case correspond à un couple $(i, j)$ d'entiers strictement positifs ($1 \leq i, j \leq n$). Pour
$n > 1$, le tableau possède $4n - 2$ \emph{diagonales} de deux types. Une diagonale du premier type est
composée des cases $(i, j)$ pour lesquelles $i + j$ est une constante, tandis qu'une diagonale du second type
est composée des cases $(i, j)$ pour lesquelles $i - j$ est une constante.

\subsection*{Problème 3}
Soit $P = A_1A_2\ldots A_k$ un polygone convexe du plan dont les sommets $A_1, A_2, \ldots, A_k$ ont des
coordonnées entières et se trouvent sur un cercle. Soit $S$ l'aire de $P$. Un entier naturel impair $n$ est
donné de sorte que le carré de la longueur de chaque côté de $P$ soit divisible par $n$.

Montrer que $2S$ est un entier divisible par $n$.

\bigskip
\begin{center}
  \emph{Second jour --- mardi 12 juillet 2016 --- Durée : 4 heures 30}
\end{center}

\subsection*{Problème 4}
Un ensemble d'entiers naturels est dit \emph{parfumé} s'il contient au moins deux éléments et si chacun de ses
éléments possède un facteur premier en commun avec au moins l'un des autres éléments. Soit
$P(n) = n^2 + n + 1$. Déterminer le plus petit entier strictement positif $b$ pour lequel il existe un entier
positif $a$ tel que l'ensemble
\[
\{P(a+1), P(a+2), \ldots, P(a+b)\}
\]
soit parfumé.

\subsection*{Problème 5}
L'équation
\[
(x-1)(x-2)\cdots(x-2016) = (x-1)(x-2)\cdots(x-2016)
\]
est écrite sur un tableau, avec $2016$ facteurs linéaires dans chaque membre. Déterminer le plus petit entier
positif $k$ pour lequel il est possible d'effacer exactement $k$ de ces $4032$ facteurs de sorte que
l'équation obtenue ait au moins un facteur dans chaque membre, mais n'ait aucune solution réelle.

\subsection*{Problème 6}
Un nombre $n \geqslant 2$ de segments se trouvent dans le plan de sorte que deux segments quelconques se
coupent en un point différent de leurs extrémités, mais sans que trois d'entre eux possèdent un point commun.
Pour chaque segment, Geoff doit choisir une extrémité et y placer une grenouille tournée vers l'autre
extrémité. Ensuite, il doit claquer des mains, au total $n - 1$ fois. Chaque fois qu'il claque des mains,
chaque grenouille saute immédiatement en avant sur son segment jusqu'au point d'intersection suivant. Les
grenouilles ne changent jamais de direction.

Geoff souhaite placer les grenouilles de sorte qu'il n'y ait jamais plus d'une grenouille à chaque point
d'intersection.
\begin{enumerate}
  \item[(a)] Montrer que si $n$ est impair, Geoff peut toujours trouver un placement des grenouilles qui
  satisfasse son souhait.
  \item[(b)] Montrer que si $n$ est pair, aucun placement des grenouilles ne satisfait son souhait.
\end{enumerate}
